\documentclass[conference]{IEEEtran}
\IEEEoverridecommandlockouts

\usepackage{cite}
\usepackage{longtable}
\usepackage{booktabs}
\usepackage{array}
\usepackage{float}
\usepackage{graphicx}
\usepackage{textcomp}
\usepackage{xcolor}
\usepackage{amsmath,amssymb,amsfonts}
\usepackage{algorithm}
\usepackage{algpseudocode}
\usepackage{hyperref}

\def\BibTeX{{\rm B\kern-.05em{\sc i\kern-.025em b}\kern-.08em
    T\kern-.1667em\lower.7ex\hbox{E}\kern-.125emX}}

\newcommand{\tightlist}{\setlength{\itemsep}{0pt}\setlength{\parskip}{0pt}}

\begin{document}

\title{Single-Machine Optimization of Perturbation-Based Loss-Grid Computation\\
\thanks{This thesis was prepared in partial fulfilment of the requirements for the Degree of Bachelor of Computer Science, Maastricht University. Supervisor(s): Max Sondag, Georgios Stamoulis}
}

\author{\IEEEauthorblockN{Given Name Surname}
\IEEEauthorblockA{\textit{Department of Advanced Computing Sciences} \\
\textit{Faculty of Science and Engineering}\\
\textit{Maastricht University}\\
Maastricht, The Netherlands}
}

\maketitle

\begin{abstract}
TODO left for last
\end{abstract}


\section{INTRODUCTION}\label{introduction}

\subsection{Motivation and Scope}\label{motivation-and-scope}

A loss landscape is a 2D
scalar field made by perturbing a trained parameter vector
along two directions and evaluating the loss at each
point~\cite{li2018losslandscape}, supporting
generalization analysis and interactive
checkpoint-comparison tools~\cite{xie2025losslens,losslens_repo}.

Computationally, the landscape is represented as an array of
scalar loss evaluations indexed by grid coordinates. Each entry
evaluates the model at the corresponding perturbed parameter
vector on data. One landscape can require hundreds to thousands
of such
evaluations~\cite{li2018losslandscape,goldstein_loss_landscape_repo}.
In interactive checkpoint comparison, this repeated work can
dominate latency on commodity hardware.

\subsection{Problem Statement}\label{problem-statement}

The problem is to evaluate candidate single-machine optimization
mechanisms for perturbation-based loss-grid computation under the
constraint of one GPU, the available CPU cores, and finite RAM
and VRAM. The optimization target is the construction of the
loss grid itself: the perturbed model evaluations that fill the
grid. The thesis studies this target through two mechanisms:
faster grid-point evaluation on one device, and CPU/GPU
distribution on one machine. Because both mechanisms depend on
hardware and workload, the study determines where each produces
statistically supported runtime reductions, where it regresses or
becomes impractical, and whether it preserves the numerical
equivalence of the resulting loss surface.

\subsection{Research Questions}\label{research-questions}

\texttt{RQ1.} Do PyTorch's stateless evaluation APIs reduce
perturbation-based loss-grid runtime, and for which architecture types?

\texttt{RQ2.} Under what conditions do idle CPU cores reduce total
loss-grid time, and how is the CPU/GPU inference throughput ratio
associated with those conditions?

\texttt{RQ3.} Does reusing the hybrid scheduler's calibrated
policy across a related checkpoint session amortize the
upfront calibration cost?

\subsection{Structure of the Report}\label{structure-of-the-report}

Section~\ref{related-work} frames the thesis with respect to prior work.
Section~\ref{methods} defines the baseline and the three
optimization approaches (one per RQ).
Section~\ref{experiments} gives the experimental design,
workloads, platforms, and measurement protocol.
Section~\ref{results} reports measurements;
Section~\ref{discussion} interprets them as applicability
conditions; Section~\ref{conclusions} answers the RQs.

\subsection{Contributions}\label{contributions}

(RQ1)~Whether the \texttt{functional\_call} + \texttt{vmap}
pattern from PyTorch's model-ensembling
tutorial~\cite{pytorch_ensembling} transfers to loss-grid
construction --- a workload with the same outer-loop shape but
different semantics (one model, many perturbations) --- across
architectures and backends.
(RQ2)~The conditions under which idle CPU cores reduce grid
time, tied to the CPU/GPU inference throughput ratio that
Gallet and Gowanlock~\cite{gallet2021heterogeneous} identified
as the core variable of heterogeneous CPU/GPU scheduling
applicability for a different workload (set operations) and
that this thesis tests in the ML-inference regime.
(RQ3)~Whether reusing the hybrid scheduler's calibrated policy
across $N=4$ related checkpoints, motivated by LossLens-style
checkpoint comparison, amortizes the upfront calibration cost
--- measured as net cumulative session speedup, with the
break-even session size reported as a diagnostic. Cases
favoring the baseline are reported as primary findings
alongside the optimization wins.


\section{RELATED WORK}\label{related-work}

Single-machine loss-grid optimization intersects loss-landscape
visualization, adaptive loss-grid sampling, heterogeneous
CPU/GPU scheduling, and autotune-and-cache patterns.

\subsection{Loss-landscape visualization and
analytics}\label{rw-losslandscape}

Goodfellow et al.~\cite{goodfellow2015qualitative} introduced
1D linear interpolation between parameter points to
qualitatively characterize neural-network optimization
trajectories. Li et al.~\cite{li2018losslandscape} extended this
to a 2D loss grid with filter-normalized perturbation directions,
the algorithm taken here as the baseline; the official
implementation~\cite{goldstein_loss_landscape_repo} runs it under
MPI on a multi-node cluster. LossLens~\cite{xie2025losslens,losslens_repo}
integrated this construct into an interactive analytics
workflow over model modes/checkpoints.

\subsection{Adaptive and progressive loss-grid
sampling}\label{rw-adaptive-sampling}

Another line of work replaces the canonical full grid with
an adaptive sampler that evaluates loss only at points predicted
to carry visual information. Prakash et al.
apply this to loss-landscape visualization directly, replacing
regular $\alpha,\beta$ sampling with an adaptive algorithm
focused on ``interesting'' regions and comparing it against a
GradVis-style baseline and a vectorized uniform
sampler~\cite{prakash2024adaptive}. Their reported outcome is
mixed: adaptive sampling beats the GradVis-style baseline but
loses to the vectorized uniform sampler due to selection
overhead, and their user-preference comparison is
inconclusive~\cite{prakash2024adaptive}.

\subsection{Heterogeneous compute and ML inference
throughput}\label{rw-scheduling}

Task-parallel heterogeneous scheduling assigns whole
independent tasks across devices. Dynamic self-scheduling,
introduced by Polychronopoulos and
Kuck~\cite{polychronopoulos1987gss}, is the simplest variant:
idle workers poll a shared queue and claim one task at a time,
requiring no per-device performance model. Richer policies in
the same family include HEFT-style cost-model scheduling and
priority-based work-stealing, exemplified by
StarPU~\cite{augonnet2011starpu}; these address load
imbalance and per-task cost variance. Gallet and
Gowanlock~\cite{gallet2021heterogeneous} identify the
CPU/GPU throughput ratio as an applicability variable for
heterogeneous task scheduling, demonstrated on a CPU--GPU
epsilon grid join: static and dynamic partitioning yield a
speedup over GPU-only execution only above a workload-dependent
ratio threshold, and regress below it.

An orthogonal approach splits a single computation across
devices, with per-device stages exchanging intermediate data
along a pipeline. OpenVINO's HETERO execution mode is a
production example at the intersection of heterogeneous
compute and ML inference~\cite{openvino_hetero}: with
pipeline-parallel \texttt{model\_distribution\_policy}, the
inference model is divided into stages, each stage is assigned
to a device, and the output of one stage is fed as input to
the next.

Production ML inference systems also pursue throughput
through request-level batching. OpenVINO's automatic
batching~\cite{openvino_auto_batching,openvino_auto_batching_src}
groups concurrent batch-1 inference requests at runtime to
improve device utilization on supported devices and
throughput-oriented modes. At implementation level,
auto-batching operates on the input/request batch dimension
of one fixed compiled model: requests are packed into slices
of a larger input/output tensor and executed by a reshaped
batched model.

\subsection{Autotune-and-cache patterns}\label{rw-autotuning}

ATLAS~\cite{whaley1998atlas} introduces an autotune-and-cache
pattern: empirically tune the inner loops of BLAS routines
(GEMM, DAXPY, etc.) for the target machine's cache and
register hierarchy; cache the selected configuration by
(machine, BLAS-op) signature; reuse it across every
subsequent matching call. The pattern is four-part: tune
once per machine, cache by signature, reuse across matching
invocations, per-machine scoping. OpenTuner~\cite{ansel2014opentuner}
canonized bounded autotuning search with explicit stopping
criteria as the discipline-level departure from ATLAS-style
exhaustive search.

\subsection{Where this thesis sits}\label{rw-placement}

Prior work provides the core
algorithm~\cite{li2018losslandscape,goldstein_loss_landscape_repo}
and its interactive checkpoint-comparison use
case~\cite{xie2025losslens,losslens_repo}, the dynamic
self-scheduling
primitive~\cite{polychronopoulos1987gss} and the broader
heterocompute-scheduler family~\cite{augonnet2011starpu},
the throughput-ratio applicability
variable~\cite{gallet2021heterogeneous} from a different
workload, the production model-parallel and request-batching
counterparts in ML
inference~\cite{openvino_hetero,openvino_auto_batching}, the
autotune-and-cache pattern~\cite{whaley1998atlas} with its
bounded-search discipline~\cite{ansel2014opentuner}, and the
adjacent adaptive-sampling
direction~\cite{prakash2024adaptive}.

The gap is the combination of workload, deployment target, and
optimization question. Existing loss-landscape work defines the
canonical grid and motivates interactive checkpoint comparison,
but does not evaluate commodity single-machine optimization of
the grid-construction runtime. Adaptive samplers change which
grid points are evaluated; this thesis keeps the canonical grid
fixed. Heterogeneous task-scheduling work supplies the
CPU/GPU-throughput-ratio variable, but from a different workload
with different task granularity and kernels. Production ML
inference systems are closer in domain, but optimize adjacent
axes: HETERO partitions one compiled model's operation graph
across devices, automatic batching groups concurrent inference
requests along the input/request batch dimension, and PyTorch's
\texttt{functional\_call}+\texttt{vmap} pattern vectorizes a
stack of parameter states on one device. None of these directly
targets the loss-grid workload shape: one model family, many
perturbations of one checkpoint, shared data, and finite CPU,
GPU, RAM, and VRAM on a commodity machine.



\section{METHODS}\label{methods}

WIP

\section{EXPERIMENTAL DESIGN}\label{experiments}

WIP

\section{RESULTS}\label{results}

WIP

\section{DISCUSSION}\label{discussion}

WIP

\section{CONCLUSIONS}\label{conclusions}

WIP

\section{REFERENCES}\label{references}

\begin{thebibliography}{00}

\bibitem{goodfellow2015qualitative}
I. J. Goodfellow, O. Vinyals, and A. M. Saxe, Qualitatively
Characterizing Neural Network Optimization Problems.
International Conference on Learning Representations (ICLR),
2015. arXiv:1412.6544.
\url{https://arxiv.org/abs/1412.6544}

\bibitem{li2018losslandscape}
H. Li, Z. Xu, G. Taylor, C. Studer, and T. Goldstein,
Visualizing the Loss Landscape of Neural Nets. Advances in
Neural Information Processing Systems 31 (NeurIPS), 2018.
arXiv:1712.09913.
\url{https://papers.nips.cc/paper/7875-visualizing-the-loss-landscape-of-neural-nets}

\bibitem{xie2025losslens}
T. Xie, J. Chen, Y. Yang, C. Geniesse, G. Shi, A. J. Chaudhari, J. K.
Cava, M. W. Mahoney, T. Perciano, G. H. Weber, and R. Maciejewski,
LossLens: Diagnostics for Machine Learning Through Loss Landscape
Visual Analytics. IEEE Computer Graphics and Applications, 2025.
arXiv:2412.13321.
\url{https://doi.org/10.1109/MCG.2024.3509374}

\bibitem{losslens_repo}
T. Xie et al., LossLens: A Visual Analytics System for Diagnosing
Machine Learning Models via Loss Landscape Analysis. GitHub
repository, accessed May 22, 2026.
\url{https://github.com/VADERASU/loss-lens-CGA}

\bibitem{prakash2024adaptive}
A. Prakash, K. Wang, and R. Balestriero,
Visualizing Loss Landscapes with Adaptive Sampling and
Gradient Free Slicing. Brown University CS 237 Course
Proceedings, 2024.
\url{https://cs.brown.edu/courses/cs237/2025/proceedings2024.pdf}

\bibitem{polychronopoulos1987gss}
C. D. Polychronopoulos and D. J. Kuck, Guided Self-Scheduling: A
Practical Scheduling Scheme for Parallel Supercomputers. IEEE
Transactions on Computers, C-36(12), 1425--1439, 1987.
\url{https://doi.org/10.1109/TC.1987.5009495}

\bibitem{gallet2021heterogeneous}
B. Gallet and M. Gowanlock, Heterogeneous CPU-GPU Epsilon Grid
Joins: Static and Dynamic Work Partitioning Strategies. Data Science
and Engineering, 2021.
\url{https://doi.org/10.1007/s41019-020-00145-x}

\bibitem{khairy2020accelsim}
M. Khairy, Z. Shen, T. M. Aamodt, and T. Rogers, Accel-Sim: An
Extensible Simulation Framework for Validated GPU Modeling. Proceedings
of the 47th Annual International Symposium on Computer Architecture
(ISCA), 2020. \url{https://doi.org/10.1109/ISCA45697.2020.00047}

\bibitem{avalos2021pka}
C. Avalos, M. Khairy, R. Green, M. Payer, and T. Rogers,
Principal Kernel Analysis: A Tractable Methodology to Simulate
Scaled GPU Workloads. Proceedings of the 54th Annual IEEE/ACM
International Symposium on Microarchitecture (MICRO), 2021.
\url{https://doi.org/10.1145/3466752.3480100}

\bibitem{gabbay2022dnnmemory}
F. Gabbay, R. Lev Aharoni, and O. Schweitzer, Deep Neural Network
Memory Performance and Throughput Modeling and Simulation Framework.
Mathematics, 2022. \url{https://doi.org/10.3390/math10214144}

\bibitem{nvidia_cuda_guide}
NVIDIA, CUDA C++ Programming Guide, section 1.1, ``The Benefits
of Using GPUs''. NVIDIA documentation, accessed May 21, 2026.
\url{https://docs.nvidia.com/cuda/archive/11.1.1/cuda-c-programming-guide/index.html}

\bibitem{openvino_hetero}
Intel, Heterogeneous Execution (HETERO). OpenVINO documentation,
accessed May 22, 2026.
\url{https://docs.openvino.ai/2024/openvino-workflow/running-inference/inference-devices-and-modes/hetero-execution.html}

\bibitem{goldstein_loss_landscape_repo}
T. Goldstein et al. loss-landscape: Code for visualizing the loss
landscape of neural nets. GitHub repository, accessed May 21, 2026.
\url{https://github.com/tomgoldstein/loss-landscape}

\bibitem{pytorch_functional_call}
PyTorch, torch.func.functional\_call API documentation. PyTorch
documentation, accessed May 21, 2026.
\url{https://docs.pytorch.org/docs/stable/generated/torch.func.functional_call.html}

\bibitem{pytorch_vmap}
PyTorch, torch.vmap API documentation. PyTorch documentation,
accessed May 21, 2026.
\url{https://docs.pytorch.org/docs/stable/generated/torch.vmap.html}

\bibitem{pytorch_ensembling}
PyTorch, Model ensembling tutorial. PyTorch documentation,
accessed May 21, 2026.
\url{https://docs.pytorch.org/tutorials/intermediate/ensembling.html}

\bibitem{whaley1998atlas}
R. C. Whaley and J. J. Dongarra, Automatically Tuned Linear
Algebra Software. Proceedings of the 1998 ACM/IEEE Conference
on Supercomputing (SC98), 1998. (Best Paper Award.)
\url{https://doi.org/10.1109/SC.1998.10004}

\end{thebibliography}

\section{APPENDICES (TODO)}\label{appendices}



\end{document}
